\chapter{Bibliography}

\section{Books}

\subsection{Common Lisp\index{Common Lisp}: the Language, 2nd Edition}

Guy L. Steele Jr., Digital Press, 1990, ISBN 1-55558-041-6

This is the book which initially defined the Common Lisp standard, before it was accepted
by ANSI. Known affectionately as ``CLtL2,'' it remains the most comprehensive reference
for ANSI CL. An on-line HTML version is available at:

\begin{center}
\texttt{http://www.cs.cmu.edu/Groups/AI/html/cltl/cltl2.html}
\end{center}

\subsection{ANSI Common Lisp\index{ANSI Common Lisp}}

Paul Graham, Prentice Hall, 1996, ISBN 0-13-370875-6

This is an excellent tutorial introduction to Common Lisp, covering all the important features
of the language in a clear and practical style. Highly recommended for beginners.

\subsection{Practical Common Lisp}

Peter Seibel, Apress, 2005, ISBN 1-59059-239-5

A hands-on introduction to Common Lisp that focuses on practical programming techniques
and real-world applications. Includes several complete example programs. Also available
online at:

\begin{center}
\texttt{http://www.gigamonkeys.com/book/}
\end{center}

\subsection{Paradigms of Artificial Intelligence Programming: Case Studies in Common Lisp}

Peter Norvig, Morgan Kaufmann, 1991, ISBN 1-55860-191-0

Known as ``PAIP,'' this book teaches both AI programming techniques and advanced Common
Lisp programming through the development of significant AI systems. Essential reading for
serious Common Lisp programmers.

\subsection{Object-Oriented Programming in Common Lisp: A Programmer's Guide to CLOS\index{CLOS}}

Sonya E. Keene, Addison-Wesley, 1989, ISBN 0-201-17589-4

The definitive guide to the Common Lisp Object System\index{Common Lisp Object System} (CLOS). Covers classes, methods,
generic functions, and the meta-object protocol.

\subsection{On Lisp: Advanced Techniques for Common Lisp}

Paul Graham, Prentice Hall, 1993, ISBN 0-13-030552-9

An advanced book focusing on macros and other sophisticated Common Lisp programming
techniques. Available online at:

\begin{center}
\texttt{http://www.paulgraham.com/onlisp.html}
\end{center}

\subsection{Common Lisp Recipes}

Edmund Weitz, Apress, 2016, ISBN 978-1-4842-1177-9

A cookbook of Common Lisp programming solutions, covering a wide range of practical
problems and techniques.

\section{Online Resources}

\subsection{Common Lisp HyperSpec}

The ANSI Common Lisp specification in hypertext format, fully cross-referenced and searchable. The authoritative reference for the language:

\begin{center}
\texttt{http://www.lispworks.com/documentation/HyperSpec/}
\end{center}

\subsection{The Common Lisp Cookbook}

A collaborative project providing recipes and techniques for common programming tasks in
Common Lisp:

\begin{center}
\texttt{http://cl-cookbook.sourceforge.net/}
\end{center}

\subsection{Planet Lisp}

An aggregator for Lisp-related blog posts and news:

\begin{center}
\texttt{http://planet.lisp.org/}
\end{center}

\subsection{Quicklisp}

A library manager for Common Lisp, making it easy to install and manage hundreds of
libraries:

\begin{center}
\texttt{http://www.quicklisp.org/}
\end{center}

\subsection{The Association of Lisp Users\index{Association of Lisp Users}}

The ALU maintains resources for the Lisp community, including a directory of Lisp systems:

\begin{center}
\texttt{http://www.alu.org/}
\end{center}

\section{Tutorials}

\subsection{Lisp Primer from Tulane University}

An interactive online tutorial for beginners:

\begin{center}
\texttt{http://www.cs.tulane.edu/www/Hennessy/courses/}

\texttt{cmps360/lisp\_primer.html}
\end{center}

\subsection{Common Lisp: A Gentle Introduction to Symbolic Computation}

David S. Touretzky, available online:

\begin{center}
\texttt{http://www.cs.cmu.edu/\~{}dst/LispBook/}
\end{center}

\subsection{Successful Lisp}

David B. Lamkins, available online:

\begin{center}
\texttt{http://successful-lisp.blogspot.com/}
\end{center}

\section{Implementation-Specific Documentation}

\subsection{Allegro CL\index{Allegro CL} Documentation}

Complete documentation for Franz Inc.'s Allegro Common Lisp:

\begin{center}
\texttt{http://www.franz.com/support/documentation/}
\end{center}

\subsection{SBCL Manual}

Documentation for Steel Bank Common Lisp:

\begin{center}
\texttt{http://www.sbcl.org/manual/}
\end{center}

\subsection{CCL Documentation}

Documentation for Clozure Common Lisp:

\begin{center}
\texttt{http://ccl.clozure.com/manual/}
\end{center}

\section{Community and Discussion}

\subsection{comp.lang.lisp}

The Usenet newsgroup for Lisp discussion, accessible via Google Groups or news readers.

\subsection{Reddit /r/lisp}

An active community for Lisp discussion:

\begin{center}
\texttt{http://www.reddit.com/r/lisp/}
\end{center}

\subsection{Lisp Discord and IRC Channels}

Active real-time chat communities for Lisp programmers can be found on Discord and
IRC (\#lisp on Libera.Chat).
